% 依赖项（必须放在主文档导言区，即 \begin{document} 之前）：
% \usepackage{tabularray}
% \UseTblrLibrary{varwidth}
% \usepackage{enumitem}
\begingroup
% \nopagebreak[4]
\enlargethispage{1cm}
\footnotesize
\centering

\DefTblrTemplate{conthead-text}{default}{（续表）}
\DefTblrTemplate{contfoot-text}{default}{续下页}

\begin{longtblr}[
  caption = {我国AI智能体领域出口管制政策汇总},
  label   = {tab:china_ai_export_control},
]{
  width   = \linewidth,
  colspec = {
    |Q[wd=0.8cm,l,m]
    |Q[wd=2.5cm,l,m]
    |Q[wd=2.5cm,l,m]
    |Q[wd=2cm,l,m]
    |X[1.00,l,m]|
  },
  cells   = {l,m},
  colsep  = 3pt,
  row{1}  = {font=\bfseries},
  rowhead = 1,
  rowsep  = 5pt,
  measure = vbox,
  stretch = -1,
}
\hline
国家 & 主管部门 & 政策法规/管制工具 & 发布/生效日期 & 核心内容 \\
\hline
\SetCell[r=2]{l,m} 中国
& \begin{itemize}[leftmargin=*,topsep=0pt,partopsep=0pt,itemsep=1pt,parsep=0pt]
    \item 商务部（牵头）
    \item 海关总署（深度参与）
\end{itemize}
& 《出口管制法》《两用物项出口管制条例》 ~\cite{中国出口管制法,两用物项出口管制条例}
& 2020年12月1日起施行；2024年12月1日起施行
& \begin{itemize}[leftmargin=*,topsep=0pt,partopsep=0pt,itemsep=1pt,parsep=0pt]
    \item 建立管制清单、临时管制、许可审批、最终用户和最终用途管理等出口管制基本制度；
    \item 两用物项出口需按清单、技术参数、目的地、最终用户、最终用途和转出口风险判断许可义务；
    \item AI相关硬件、软件、技术或模型只有在落入清单、临时管制或相关最终用途规则时，才纳入出口管制。
\end{itemize} \\
\cline{2-5}
& \begin{itemize}[leftmargin=*,topsep=0pt,partopsep=0pt,itemsep=1pt,parsep=0pt]
    \item 商务部（牵头）
    \item 科技部（深度参与）
\end{itemize}
& 《中国禁止出口限制出口技术目录》~\cite{中国禁止出口限制出口技术目录}
& 2023.12发布，2025.07调整
& \begin{itemize}[leftmargin=*,topsep=0pt,partopsep=0pt,itemsep=1pt,parsep=0pt]
    \item 将人工智能交互界面、语音合成、个性化信息推送等信息处理技术纳入限制出口技术范围；
    \item 属于限制出口目录的技术，应当依照技术出口许可程序办理相关手续；
    \item 目录所列技术同时属于军民两用技术的，纳入两用物项出口管制体系管理。
\end{itemize} \\
\hline
\end{longtblr}
\endgroup
